\chapter{组合数学}

\section{加法原理与乘法原理}

\subsection{加法原理}

在生活中，完成一件事通常有不同的方法和步骤。如果完成一件事有 $n$ 种不同的方法，每种方法的具体方法数分别为 $a_1, a_2, \dots, a_n$，那么总的方法数就是 $\sum_{i=1}^n a_i$。这被称为\textbf{加法原理}。

\subsection{乘法原理}

如果完成一件事有 $n$ 个步骤，每个步骤的方法数分别为 $a_1, a_2, \cdot, a_n$，那么总的方法数就是 $\prod_{i=1}^n a_i$。这被称为\textbf{乘法原理}。

\section{排列数与组合数}

\subsection{排列数}

从 $n$ 个元素中任取 $m$ 个元素，按一定的顺序排列，则总的方案数记作 $\operatorname{A}_n^m$。$\operatorname{A}_n^m$ 的计算方式如下：

\[
\operatorname{A}_n^m = n(n-1)(n-2) \cdots (n-m+1) = \frac{n!}{(n-m)!}
\]

\subsection{组合数}

从 $n$ 个元素中任取 $m$ 个元素，不考虑顺序排列，则总的方案数记作 $\binom{n}{m}$ 或 $\operatorname{C}_n^m$\hspace{0.3em}\footnote{\quad 本资料除本处和附录 A 外统一使用 $\operatorname{C}_n^m$，主要是为了考虑高中生。更规范的表示方法是 $\binom{n}{m}$。}。$\operatorname{C}_n^m$ 的计算方式如下：

\[
\operatorname{C}_n^m = \frac{\operatorname{A}_n^m}{\operatorname{A}_m^m} = \frac{n!}{m!(n-m)!}
\]

特别地，当 $m > n$ 时，$\operatorname{A}_n^m = \operatorname{C}_n^m = 0$。特别地，$\operatorname{C}_n^0 = 1$。

由于从 $n$ 中选出 $m$ 个物品和从 $n$ 个物品中不选出 $(n-m)$ 个物品是一样的，$\operatorname{C}_n^m = \operatorname{C}_n^{n-m}$。 

\section{杨辉三角与二项式定理}

\subsection{杨辉三角}

有一个三角形，内部有若干个点：第 $1$ 行有 $1$ 个点，第 $2$ 行有 $2$ 个点，以此类推。第 $1$ 行的数为 $1$，随后每行的每个数均为其左上方和右上方数的和（若无则视为 $0$），就可以构成一个数字三角形。这样的数字三角形叫作\textbf{杨辉三角}（又称贾宪三角，外国称为帕斯卡三角形）。

\begin{center}
$1$\\
$1$ \quad $1$\\
$1$ \quad $2$ \quad $1$\\
$1$ \quad $3$ \quad $3$ \quad $1$\\
$1$ \quad $4$ \quad $6$ \quad $4$ \quad $1$\\
$\cdots\cdots$
\end{center}

杨辉三角有一个与组合数强相关的性质：除第 $1$ 行外，第 $i+1$ 行从左往右数第 $j+1$ 个数恰为 $\operatorname{C}_i^j$。由此，我们可以得到组合数的递推算法：

\[
\operatorname{C}_n^m = \operatorname{C}_{n-1}^{m-1} + \operatorname{C}_{n-1}^m
\]

这个等式叫作帕斯卡恒等式。

\subsection{二项式定理}

我们知道，$a+b = a+b$，$(a+b)^2 = a^2 + 2ab + b^2$，我们注意到其各项系数恰为 $1$、$1$ 和 $1$、$2$、$1$，刚好是杨辉三角的第 $2$ 行和第 $3$ 行。更高次数的展开有没有这种规律呢？

经过探究，数学家门发现了如下的规律：对于一个二项式 $(a+b)$ 的 $n$ 次幂的展开式，若按从左往右 $a$ 的次数逐渐减小，$b$ 的次数逐渐增大的顺序排列，那么各项的系数刚好是杨辉三角的第 $n+1$ 行。这被称为\textbf{二项式定理}。又由于杨辉三角和组合数的关系，因此，二项式定理可以表示成如下形式：

\[
(a+b)^n = \sum_{i=0}^n \operatorname{C}_n^i a^{n-i} b^i
\]